% !TEX TS-program = lualatex
% encoding : utf8
% Copyright 2011-2026  Alain Matthes
% SPDX-License-Identifier: LPPL-1.3c
% Maintainer: Alain Matthes

\PassOptionsToPackage{unicode}{hyperref}

\documentclass[DIV         = 14,
               fontsize    = 10,
               index       = totoc,
               twoside,
               cadre,
               headings    = small,
               ]{tkz-doc}
\gdef\tkznameofpack{tkz-euclide}
\gdef\tkzversionofpack{5.15c}
\gdef\tkzdateofpack{\today}
\gdef\tkznameofdoc{doc-tkz-euclide}
\gdef\tkzversionofdoc{5.15c}
\gdef\tkzdateofdoc{\today}
\gdef\tkzauthorofpack{Alain Matthes}
\gdef\tkzadressofauthor{}
\gdef\tkznamecollection{AlterMundus}
\gdef\tkzurlauthor{http://altermundus.fr}
\gdef\tkzengine{lualatex}
\gdef\tkzurlauthorcom{http://altermundus.fr}
\nameoffile{\tkznameofpack}
% -- Packages ---------------------------------------------------
\usepackage[dvipsnames,svgnames]{xcolor}
\usepackage{calc}
\usepackage{tkz-base}
\usepackage[lua]{tkz-euclide}
\usepackage{pgfornament}
\usetikzlibrary{backgrounds}
\usepackage[colorlinks,pdfencoding=auto, psdextra]{hyperref}
\hypersetup{
      linkcolor=Gray,
      citecolor=Green,
      filecolor=Mulberry,
      urlcolor=NavyBlue,
      menucolor=Gray,
      runcolor=Mulberry,
      linkbordercolor=Gray,
      citebordercolor=Green,
      filebordercolor=Mulberry,
      urlbordercolor=NavyBlue,
      menubordercolor=Gray,
      runbordercolor=Mulberry,
      pdfsubject={Euclidean Geometry},
      pdfauthor={\tkzauthorofpack},
      pdftitle={\tkznameofpack},
      pdfcreator={\tkzengine},
      bookmarksdepth=3
}
\usepackage{tkzexample}
\usepackage{fontspec}
\setmainfont{texgyrepagella}[
  Extension = .otf,
  UprightFont = *-regular ,
  ItalicFont  = *-italic  ,
  BoldFont    = *-bold    ,
  BoldItalicFont = *-bolditalic
]
\setsansfont{texgyreheros}[
  Extension = .otf,
  UprightFont = *-regular ,
  ItalicFont  = *-italic  ,
  BoldFont    = *-bold    ,
  BoldItalicFont = *-bolditalic ,
]

\setmonofont{lmmono10-regular.otf}[
  Numbers={Lining,SlashedZero},
  ItalicFont=lmmonoslant10-regular.otf,
  BoldFont=lmmonolt10-bold.otf,
  BoldItalicFont=lmmonolt10-boldoblique.otf,
]
\newfontfamily\ttcondensed{lmmonoltcond10-regular.otf}
%% (La)TeX font-related declarations:
\linespread{1.05}      % Pagella needs more space between lines
\usepackage[math-style=literal,bold-style=literal]{unicode-math}
\usepackage{fourier-otf}
\let\rmfamily\ttfamily
\usepackage{multicol,lscape}
\usepackage[english]{babel}
\usepackage[normalem]{ulem}
\usepackage{multirow,multido,booktabs,cellspace}
\usepackage{shortvrb,fancyvrb,bookmark}
\usepackage{makeidx}
\makeindex

%<--------------------------------------------------------------------------->
% settings styles
\tkzSetUpColors[background=white,text=black]
\tkzSetUpCompass[color=orange, line width=.2pt,delta=10]
\tkzSetUpArc[color=gray,line width=.2pt,delta=10]
\tkzSetUpPoint[size=2,color=teal]
\tkzSetUpLine[line width=.2pt,color=teal]
\tkzSetUpStyle[color=orange,line width=.2pt]{new}
\tikzset{every picture/.style={line width=.2pt}}
\tikzset{label angle style/.append style={color=teal,font=\footnotesize}}
\tikzset{label style/.append style={below,color=teal,font=\scriptsize}}
\tikzset{new/.style={color=orange,line width=.2pt}}

\AtBeginDocument{\MakeShortVerb{\|}} % link to shortvrb
\setlength{\emergencystretch}{3em}% fewer overfull lines with long macro names
\def\code{\texttt}
\newcommand*{\ItkzPopt}[2]{\texttt{#2}\index{#1_3@\texttt{#1: options}!\texttt{#2}}}

\begin{document}

\parindent=0pt
\tkzTitleFrame{tkz-euclide v5\\Euclidean Geometry}
\clearpage

\defoffile{\lefthand\ \tkzname{\tkznameofpack} is a set of convenient macros for drawing figures in the plane with a Cartesian coordinate system. It handles the most classic situations of Euclidean geometry. \tkzname{\tkznameofpack} is built on top of PGF and its associated front-end \TIKZ\ and is a (La)TeX-friendly drawing package. It provides a high-level interface to build figures with little effort: you construct them step by step, as closely as possible to a construction made by hand.\\
Version 5 of \tkzname{\tkznameofpack} can use Lua for some calculations; see the sections \code{News} and \code{Lua} at the end of this documentation.\\
Please note: English is not my native language, so there may be some errors.
}

\presentation

\vspace*{1cm}
\lefthand\ First of all, I would like to thank \textbf{Till Tantau} for the beautiful package \href{https://github.com/pgf-tikz/pgf}{\TIKZ}.

\vspace*{12pt}
\lefthand\ Acknowledgements: I received much valuable advice, remarks, corrections and examples from

\tkzimp{Jean-Côme Charpentier}, \tkzimp{Josselin Noirel}, \tkzimp{Manuel Pégourié-Gonnard}, \tkzimp{Franck Pastor}, \tkzimp{David Arnold},

\tkzimp{Ulrike Fischer}, \tkzimp{Stefan Kottwitz}, \tkzimp{Christian Tellechea}, \tkzimp{Nicolas Kisselhoff}, \tkzimp{Wolfgang Büchel},

\tkzimp{John Kitzmiller}, \tkzimp{Dimitri Kapetas}, \tkzimp{Gaétan Marris}, \tkzimp{Mark Wibrow}, \tkzimp{Yves Combe}, \tkzimp{Paul Gaborit}, \tkzimp{Laurent Van Deik} and \tkzimp{Muzimuzhi Z}.

\vspace*{12pt}
\lefthand\ I would also like to thank Eric Weisstein, creator of MathWorld:
\href{http://mathworld.wolfram.com/about/author.html}{MathWorld}.

\vspace*{12pt}
\lefthand\ You can find some examples on my site:
\href{http://altermundus.fr}{altermundus.fr}. \hspace{2cm} under construction!

\vfill
Please report typos or any other comments about this documentation to \href{mailto:al.ma@mac.com}{\textcolor{blue}{Alain Matthes}}.

This file can be redistributed and/or modified under the terms of the \LaTeX{}
Project Public License, distributed from \href{http://www.ctan.org/}{CTAN}\  archives.

\clearpage
\setcounter{tocdepth}{2}% sections and subsections; subsubsections stay in the PDF bookmarks
\tableofcontents
\clearpage
\newpage

\part{Getting started}
\input{TKZdoc-euclide-presentation.tex}
\input{TKZdoc-euclide-installation.tex}
\input{TKZdoc-euclide-elements.tex}
\input{TKZdoc-euclide-documentation.tex}

\part{Points}
\input{TKZdoc-euclide-points.tex}
\input{TKZdoc-euclide-pointsSpc.tex}
\input{TKZdoc-euclide-pointby.tex}
\input{TKZdoc-euclide-pointwith.tex}
\input{TKZdoc-euclide-rnd.tex}

\part{Lines, triangles and polygons}
\input{TKZdoc-euclide-lines.tex}
\input{TKZdoc-euclide-triangles.tex}
\input{TKZdoc-euclide-polygons.tex}

\part{Circles}
\newpage
\section{Circles}

Among the following macros, one draws a circle, which is not a real feat: you need the center of the circle and either its radius or a point of the circle. The most frequent use seemed to me to be a circle with a given center through a given point; this is the default, otherwise you use the option \tkzname{R}. There are many special circles, for example the circumcircle of a triangle.

\begin{itemize}
  \item a first macro, \tkzcname{tkzDefCircle}, gives the center and the radius (in cm) of a particular circle, which you retrieve with the macros \tkzcname{tkzGetPoint} and \tkzcname{tkzGetLength};
 
 \item a macro to draw a circle, \tkzcname{tkzDrawCircle};
 
 \item a macro to color a disc without drawing the circle, \tkzcname{tkzFillCircle};
 
 \item sometimes a drawing must be contained in a disc: this is the role of \tkzcname{tkzClipCircle};
  
 \item finally, \tkzcname{tkzLabelCircle} gives a label to a circle.
\end{itemize} 

\subsection{Characteristics of a circle: \tkzcname{tkzDefCircle}}
 
This macro gives the characteristics (center and radius) of some circles.

\begin{NewMacroBox}{tkzDefCircle}{\oarg{local options}\parg{A,B} or \parg{A,B,C}}%
\tkzHandBomb\ Note that the arguments are lists of two or three points. The macro is used either with \\ \tkzcname{tkzGetPoints}, to get the center and a point of the circle, or directly with \\ \tkzname{tkzFirstPointResult} and \tkzname{tkzSecondPointResult} if you don't need to keep the results. \tkzcname{tkzGetLength} gives the radius.

\medskip
\begin{tabular}{lll}%
\toprule
arguments           & example & explanation                         \\
\midrule
\TAline{\parg{pt1,pt2} or \parg{pt1,pt2,pt3}}{\parg{A,B}} {$[AB]$ is radius $A$ is the center}
\bottomrule
\end{tabular} 

\medskip
\begin{tabular}{lll}%
\toprule
options             & default & definition                         \\ 
\midrule
\TOline{R}       {R}{circle characterized by a center and a radius} 
\TOline{diameter}{}{circle characterized by two points defining a diameter}
\TOline{circum}       {}{circle circumscribed of a triangle} 
\TOline{in}           {}{incircle a triangle} 
\TOline{ex}           {}{excircle of a  triangle}
\TOline{euler or nine}{}{Euler's Circle}
\TOline{spieker}      {}{Spieker Circle}
\TOline{apollonius}   {}{circle of Apollonius}
\TOline{orthogonal from} {}{[orthogonal from = A ](O,M)}
\TOline{orthogonal through}{}{[orthogonal through = A and B](O,M)} 
\TOline{K} {1}{coefficient used for a circle of Apollonius} 
 \bottomrule
\end{tabular}

\medskip
\emph{In the following examples, I draw the circles with a macro not yet presented. You may only need the center and a point of the circle.}
\end{NewMacroBox} 

\subsubsection{Example with option \tkzname{R}}  
With the macro \tkzcname{tkzGetPoint} we get a point of the circle, its "east pole".

\begin{tkzexample}[latex=7cm,small]  
\begin{tikzpicture}[scale=1]
  \tkzDefPoint(3,3){C}
  \tkzDefPoint(5,5){A}
   \tkzCalcLength(A,C) \tkzGetLength{rAC}
  \tkzDefCircle[R](C,\rAC) \tkzGetPoint{B}
  \tkzDrawCircle(C,B)
  \tkzDrawSegment(C,A)
  \tkzLabelSegment[above left](C,A){$2\sqrt{2}$}
  \tkzDrawPoints(A,B,C)
  \tkzLabelPoints(A,C,B)
\end{tikzpicture} 
\end{tkzexample}     


\subsubsection{Example with option \tkzname{diameter}}  
 It is simpler here to search directly for the midpoint of $[AB]$. The result is the center and if necessary 
\begin{tkzexample}[latex=7cm,small]  
\begin{tikzpicture}
  \tkzDefPoint(0,0){O}
  \tkzDefPoint(2,2){B}
  \tkzDefCircle[diameter](O,B) \tkzGetPoint{A}
  \tkzDrawCircle(A,B)
  \tkzDrawPoints(O,A,B)
  \tkzDrawSegment(O,B)
  \tkzLabelPoints(O,A,B)
  \tkzLabelSegment[above left](O,A){$\sqrt{2}$}
  \tkzLabelSegment[above left](A,B){$\sqrt{2}$}
  \tkzMarkSegments[mark=s||](O,A A,B)
\end{tikzpicture}
\end{tkzexample}    

\subsubsection{Circles inscribed and circumscribed for a given triangle} 
 
\begin{tkzexample}[latex=7cm,small]  
\begin{tikzpicture}[scale=.75]
 \tkzDefPoint(2,2){A}  \tkzDefPoint(5,-2){B}
 \tkzDefPoint(1,-2){C}
 \tkzDefCircle[in](A,B,C)
 \tkzGetPoints{I}{x}   
 \tkzDefCircle[circum](A,B,C)
 \tkzGetPoint{K}  
 \tkzDrawCircles[new](I,x K,A) 
 \tkzLabelPoints[below](B,C)
 \tkzLabelPoints[above left](A,I,K)
 \tkzDrawPolygon(A,B,C)
 \tkzDrawPoints(A,B,C,I,K) 
\end{tikzpicture} 
\end{tkzexample}

\subsubsection{Example with option \tkzname{ex}}
We want to define the excircle of a triangle opposite the vertex $C$.

\begin{tkzexample}[latex=8cm,small]
\begin{tikzpicture}[scale=.75]
  \tkzDefPoints{ 0/0/A,4/0/B,0.8/4/C}
  \tkzDefCircle[ex](B,C,A)                   
  \tkzGetPoints{J_c}{h}
  \tkzDefPointBy[projection=onto A--C ](J_c)   
  \tkzGetPoint{X_c}
  \tkzDefPointBy[projection=onto A--B ](J_c)   
  \tkzGetPoint{Y_c}     
  \tkzDefCircle[in](A,B,C)    
  \tkzGetPoints{I}{y}
  \tkzDrawCircles[color=lightgray](J_c,h I,y)
  \tkzDefPointBy[projection=onto A--C ](I)
  \tkzGetPoint{F}
  \tkzDefPointBy[projection=onto A--B ](I)
  \tkzGetPoint{D}
  \tkzDrawPolygon(A,B,C)
  \tkzDrawLines[add=0 and 1.5](C,A C,B)
  \tkzDrawSegments(J_c,X_c I,D  I,F J_c,Y_c)
  \tkzMarkRightAngles(A,F,I B,D,I J_c,X_c,A
     J_c,Y_c,B)
  \tkzDrawPoints(B,C,A,I,D,F,X_c,J_c,Y_c)
  \tkzLabelPoints(B,A,J_c,I,D)
  \tkzLabelPoints[above](Y_c)
  \tkzLabelPoints[left](X_c)
  \tkzLabelPoints[above left](C)
  \tkzLabelPoints[left](F)
\end{tikzpicture}  
\end{tkzexample}

\subsubsection{Euler's circle for a given triangle with option \tkzname{euler}}
 
We verify that this circle passes through the midpoint of each side.
\begin{tkzexample}[latex=6cm,small]  
\begin{tikzpicture}[scale=.75]
   \tkzDefPoint(5,3.5){A} 
   \tkzDefPoint(0,0){B} \tkzDefPoint(7,0){C}
   \tkzDefCircle[euler](A,B,C)
   \tkzGetPoints{E}{e}
   \tkzDefSpcTriangle[medial](A,B,C){M_a,M_b,M_c}
   \tkzDrawCircle[new](E,e)
   \tkzDrawPoints(A,B,C,E,M_a,M_b,M_c)    
   \tkzDrawPolygon(A,B,C)    
   \tkzLabelPoints[below](B,C)  
   \tkzLabelPoints[left](A,E)   
\end{tikzpicture}
\end{tkzexample}

\subsubsection{Apollonius circles for a given segment: option \tkzname{apollonius}} 
 
\begin{tkzexample}[latex=9cm,small]    
\begin{tikzpicture}[scale=0.75]
  \tkzDefPoint(0,0){A} 
  \tkzDefPoint(4,0){B}
  \tkzDefCircle[apollonius,K=2](A,B)
  \tkzGetPoints{K1}{x}
  \tkzDrawCircle[color = teal!50!black,
      fill=teal!20,opacity=.4](K1,x)
  \tkzDefCircle[apollonius,K=3](A,B)
  \tkzGetPoints{K2}{y}
  \tkzDrawCircle[color=orange!50,
      fill=orange!20,opacity=.4](K2,y) 
  \tkzLabelPoints[below](A,B,K1,K2)
  \tkzDrawPoints(A,B,K1,K2) 
  \tkzDrawLine[add=.2 and 1](A,B)  
\end{tikzpicture}
\end{tkzexample}  

\subsubsection{Excircles of a given triangle: option \tkzname{ex}}
 You can also get the center and its projection on one side of the triangle
 
 with \tkzcname{tkzGetFirstPoint\{Jb\}} and \tkzcname{tkzGetSecondPoint\{Tb\}}.
 
\begin{tkzexample}[latex=8cm,small]  
\begin{tikzpicture}[scale=.6]
  \tkzDefPoint(0,0){A}
  \tkzDefPoint(3,0){B}
  \tkzDefPoint(1,2.5){C}
  \tkzDefCircle[ex](A,B,C) \tkzGetPoints{I}{i}
  \tkzDefCircle[ex](C,A,B) \tkzGetPoints{J}{j}
  \tkzDefCircle[ex](B,C,A) \tkzGetPoints{K}{k}
  \tkzDefCircle[in](B,C,A) \tkzGetPoints{O}{o}
  \tkzDrawCircles[new](J,j I,i K,k O,o) 
  \tkzDrawLines[add=1.5 and 1.5](A,B A,C B,C)
  \tkzDrawPolygon[purple](I,J,K)
  \tkzDrawSegments[new](A,K B,J C,I)
  \tkzDrawPoints(A,B,C)
  \tkzDrawPoints[new](I,J,K)   
  \tkzLabelPoints(A,B,C,I,J,K)
\end{tikzpicture}
\end{tkzexample}
 
\subsubsection{Spieker circle with option \tkzname{spieker}}   
The incircle of the medial triangle $M_aM_bM_c$ is the Spieker circle:

\begin{tkzexample}[latex=6cm, small]
\begin{tikzpicture}[scale=1.25]
  \tkzDefPoints{ 0/0/A,4/0/B,0.8/4/C}
  \tkzDefSpcTriangle[medial](A,B,C){M_a,M_b,M_c}
  \tkzDefTriangleCenter[spieker](A,B,C) 
  \tkzGetPoint{S_p}
  \tkzDrawPolygon(A,B,C)
  \tkzDrawPolygon[cyan](M_a,M_b,M_c)
  \tkzDrawPoints(B,C,A)
  \tkzDefCircle[spieker](A,B,C)
  \tkzDrawPoints[new](M_a,M_b,M_c,S_p)
  \tkzDrawCircle[new](tkzFirstPointResult,%
    tkzSecondPointResult)
  \tkzLabelPoints[right](M_a)
  \tkzLabelPoints[left](M_b)
  \tkzLabelPoints[below](A,B,M_c,S_p)
  \tkzLabelPoints[above](C)
\end{tikzpicture}
\end{tkzexample}
 
\subsection{Projection of excenters}

\begin{NewMacroBox}{tkzDefProjExcenter}{\oarg{local options}\parg{A,B,C}\parg{a,b,c}\marg{X,Y,Z}}%
Each excenter has three projections on the sides of the triangle $ABC$. We can get them all with one macro:\\ \tkzcname{tkzDefProjExcenter[name=J](A,B,C)(a,b,c)\{Y,Z,X\}}.

\medskip
\begin{tabular}{lll}%
\toprule
options             & default & definition                        \\
\midrule
\TOline{name} {no defaut}{used to name the vertices}
\bottomrule
\end{tabular}

\begin{tabular}{lll}%
arguments & default & definition \\
\midrule
\TAline{(pt1=$\alpha_1$,pt2=$\alpha_2$,\dots)}{no default}{Each point has a assigned weight}
\end{tabular}

\medskip
\end{NewMacroBox}

\subsubsection{\tkzname{Excircles}}

\begin{tkzexample}[vbox,small]
\begin{tikzpicture}[scale=.6]
\tikzset{line style/.append style={line width=.2pt}}
\tikzset{label style/.append style={color=teal,font=\footnotesize}}
\tkzDefPoints{0/0/A,5/0/B,0.8/4/C}
\tkzDefSpcTriangle[excentral,name=J](A,B,C){a,b,c} 
\tkzDefSpcTriangle[intouch,name=I](A,B,C){a,b,c}
\tkzDefProjExcenter[name=J](A,B,C)(a,b,c){X,Y,Z}
\tkzDefCircle[in](A,B,C)   \tkzGetPoint{I} \tkzGetSecondPoint{T}  
\tkzDrawCircles[red](Ja,Xa Jb,Yb Jc,Zc)
\tkzDrawCircle(I,T) 
\tkzDrawPolygon[dashed,color=blue](Ja,Jb,Jc)
\tkzDrawLines[add=1.5 and 1.5](A,C A,B B,C)
\tkzDrawSegments(Ja,Xa Ja,Ya Ja,Za
                 Jb,Xb Jb,Yb Jb,Zb
                 Jc,Xc Jc,Yc Jc,Zc
                 I,Ia I,Ib I,Ic)
\tkzMarkRightAngles[size=.2,fill=gray!15](Ja,Za,B Ja,Xa,B Ja,Ya,C Jb,Yb,C)
\tkzMarkRightAngles[size=.2,fill=gray!15](Jb,Zb,B Jb,Xb,C Jc,Yc,A Jc,Zc,B)
\tkzMarkRightAngles[size=.2,fill=gray!15](Jc,Xc,C I,Ia,B I,Ib,C I,Ic,A)
\tkzDrawSegments[blue](Jc,C Ja,A Jb,B)
\tkzDrawPoints(A,B,C,Xa,Xb,Xc,Ja,Jb,Jc,Ia,Ib,Ic,Ya,Yb,Yc,Za,Zb,Zc)
\tkzLabelPoints(A,Ya,Yb,Ja,I)
\tkzLabelPoints[left](Jb,Ib,Yc)
\tkzLabelPoints[below](Zb,Ic,Jc,B,Za,Xa)
\tkzLabelPoints[above right](C,Zc,Yb)
\tkzLabelPoints[right](Xb,Ia,Xc)
\end{tikzpicture}
\end{tkzexample}
 
\subsubsection{\tkzname{Orthogonal from}}
Circle with a given center, orthogonal to a given circle. \tkzcname{tkzGetPoints\{z1\}\{z2\}} gives two points of the circle.

\begin{tkzexample}[latex=7cm,small]
\begin{tikzpicture}[scale=.75]
  \tkzDefPoints{0/0/O,1/0/A}
  \tkzDefPoints{1.5/1.25/B,-2/-3/C}
  \tkzDefCircle[orthogonal from=B](O,A)
  \tkzGetPoints{z1}{z2}
  \tkzDefCircle[orthogonal from=C](O,A)
  \tkzGetPoints{t1}{t2}
  \tkzDrawCircle(O,A)
  \tkzDrawCircles[new](B,z1 C,t1)
  \tkzDrawPoints(t1,t2,C)
  \tkzDrawPoints(z1,z2,O,A,B)
  \tkzLabelPoints[right](O,A,B,C)
\end{tikzpicture}
\end{tkzexample}

\subsubsection{\tkzname{Orthogonal through}}
Circle orthogonal to a given circle, through two given points.
\begin{tkzexample}[latex=6cm,small]
\begin{tikzpicture}[scale=1]
  \tkzDefPoint(0,0){O}
  \tkzDefPoint(1,0){A}
  \tkzDrawCircle(O,A)
  \tkzDefPoint(-1.5,-1.5){z1}
  \tkzDefPoint(1.5,-1.25){z2}
  \tkzDefCircle[orthogonal through=z1 and z2](O,A)
   \tkzGetPoint{c}
  \tkzDrawCircle[new](tkzPointResult,z1)
  \tkzDrawPoints[new](O,A,z1,z2,c)
  \tkzLabelPoints[right](O,A,z1,z2,c)
\end{tikzpicture}
\end{tkzexample}

\endinput
\input{TKZdoc-euclide-circleby.tex}

\part{Intersections and angles}
\section{Intersections}

You can determine the points of intersection of two straight lines, of a straight line and a circle, and of two circles.

The associated commands do not check that the intersection points exist: this is up to the user (see the tests \tkzcname{tkzTestInterLC} and \tkzcname{tkzTestInterCC}).

\subsection{Intersection of two straight lines \tkzcname{tkzInterLL}}
\begin{NewMacroBox}{tkzInterLL}{\parg{$A,B$}\parg{$C,D$}}%
Defines the intersection point \tkzname{tkzPointResult} of the two lines $(AB)$ and $(CD)$. The known points are given in pairs (two per line) between parentheses, and the resulting point is retrieved with the macro \tkzcname{tkzGetPoint}.
\end{NewMacroBox}

\subsubsection{Example of intersection between two straight lines}

\begin{tkzexample}[latex=7cm,small]
\begin{tikzpicture}[rotate=-45,scale=.75]
  \tkzDefPoint(2,1){A}   
  \tkzDefPoint(6,5){B}
  \tkzDefPoint(3,6){C}   
  \tkzDefPoint(5,2){D}
  \tkzDrawLines(A,B C,D)
  \tkzInterLL(A,B)(C,D)  
     \tkzGetPoint{I}
  \tkzDrawPoints[color=blue](A,B,C,D)
   \tkzDrawPoint[color=red](I)
\end{tikzpicture}
\end{tkzexample}

\subsection{Intersection of a straight line and a circle  \tkzcname{tkzInterLC}}

As before, the line is defined by a pair of points. The circle
 is also defined by a pair:
\begin{itemize}
\item $(O,C)$, a pair of points: the first one is the center and the second one is any point of the circle;
\item $(O,r)$, where $r$ is the radius.
\end{itemize}

\begin{NewMacroBox}{tkzInterLC}{\oarg{options}\parg{$A,B$}\parg{$O,C$} or \parg{$O,r$} or \parg{$O,C,D$}}%
So the arguments are two pairs.

\medskip
\begin{tabular}{lll}%
\toprule
options            & default & definition                         \\ 
\midrule
\TOline{N}         {N}    {(O,C) determines the circle}
\TOline{R}         {N}    {(O, 1 ) unit 1 cm}  
\TOline{with nodes}{N}    {(O,C,D) CD is a radius}  
\TOline{common=pt} {}     {pt is common point; tkzFirstPoint gives the other point}
\TOline{near}      {}     {tkzFirstPoint is the closest point to the first point of the line}
\bottomrule
\end{tabular}

\medskip   
The macro defines the intersection points $I$ and $J$ of the line $(AB)$ and the circle with center $O$ and radius $r$, if they exist; otherwise, an error is reported in the |.log| file. With \tkzname{with nodes}, you don't need to calculate the radius, which is the length of $[CD]$.
If \tkzname{common} and \tkzname{near} are not used, \tkzname{tkzFirstPoint} is the point that gives the smallest angle (formed with \tkzname{tkzSecondPoint} and the center of the circle).
If the line passes through the center, \tkzname{tkzSecondPoint} is the point located in the direction of the line, from $A$ to $B$ (since version 5.15).
\end{NewMacroBox}

\begin{NewMacroBox}{tkzTestInterLC}{\parg{$O,A$}\parg{$O',B$}}%
The arguments are two pairs, which define a line and a circle (with a center and a point of the circle). If the line and the circle intersect, the test \tkzcname{iftkzFlagLC} is true.
\end{NewMacroBox}

\subsubsection{Testing a line-circle intersection: \tkzcname{tkzTestInterLC}}

\begin{tkzexample}[latex=7cm,small]
  \begin{tikzpicture}[scale=1]
    \tkzDefPoints{% x   y   name
                    3    /4    /I,
                    3    /2    /P,
                    0    /2    /La,
                    8    /3    /Lb}
  \tkzDrawCircle(I,P)
  \foreach \i in {1,...,3}{%
     \coordinate  (Lb) at (8,\i);
     \tkzDrawLine(La,Lb)
     \tkzTestInterLC(La,Lb)(I,P)
      \iftkzFlagLC
      \tkzInterLC(La,Lb)(I,P)  
      \tkzGetPoints{a}{b}
      \tkzDrawPoints(a,b)
      \fi
     }
  \end{tikzpicture}
\end{tkzexample}


\subsubsection{Line-circle intersection}

In the following example, the circle is drawn with two points, and the intersection of the straight line and the circle uses two pairs of points. We will compare the angles $\widehat{D,E,O}$ and $\widehat{E,D,O}$. These angles have opposite orientations. \tkzname{tkzFirstPoint} is assigned to the point that forms the angle with the smallest measure (counterclockwise direction). The counterclockwise angle $\widehat{D,E,O}$ measures $360^\circ$ minus the measure of $\widehat{O,E,D}$.

\begin{tkzexample}[latex=7cm,small]
\begin{tikzpicture}[scale=.75]
 \tkzInit[xmax=5,ymax=4]
 \tkzDefPoint(1,1){O} 
 \tkzDefPoint(-2,4){La} 
 \tkzDefPoint(5,0){Lb} 
 \tkzDefPoint(3,3){C}
 \tkzInterLC(La,Lb)(O,C)  \tkzGetPoints{D}{E}  
 \tkzMarkAngle[->,size=1.5](E,D,O)
 \tkzDrawPolygons[new](O,D,E)
 \tkzMarkAngle[->,size=1.5](D,E,O)
 \tkzDrawCircle(O,C)
 \tkzDrawPoints[color=teal](O,La,Lb,C)
 \tkzDrawPoints[color=red](D,E)
 \tkzDrawLine(La,Lb)
 \tkzLabelPoints[above right](O,La,Lb,C,D,E)
\end{tikzpicture} 
\end{tkzexample}

\subsubsection{Line passing through the center}\label{interLC-center}
This case is special: the two points and the center are aligned, so the angles cannot be compared. Since version 5.15, the order is fixed: with \tkzcname{tkzInterLC}\parg{La,Lb}\parg{O,D}, \tkzname{tkzSecondPoint} is the point located in the direction of the line, from $La$ to $Lb$, and \tkzname{tkzFirstPoint} is the other one. (Before 5.15, the order depended on rounding, and the results with and without the option \tkzname{lua} could be different.)

You can also use the option \tkzname{near}: \tkzname{tkzFirstPoint} is then the point closest to the first point given for the line. Here we want $A$ to be closest to $Lb$; \tkzcname{tkzInterLC}\parg{Lb,La}\parg{O,D} without option gives the same point.

\begin{tkzexample}[latex=8cm,small]
\begin{tikzpicture}
\tkzDefPoints{% x   y   name
             0    /1    /D,
             6    /0    /B,
             3    /3    /O,
             2    /2    /La,
             5    /5    /Lb}
  \tkzDrawCircle(O,D)
  \tkzDrawLine(La,Lb)
  \tkzInterLC[near](Lb,La)(O,D)  
  \tkzGetFirstPoint{A}
  \tkzDrawSegments(O,A)
  \tkzDrawPoints(O,D,A,La,Lb)
  \tkzLabelPoints(O,D,A,La,Lb)
\end{tikzpicture}
\end{tkzexample}

\subsubsection{Line-circle intersection with option \tkzname{common}}
A special case that often occurs: a point of the line is on the circle, and we look for the other common point.
\begin{tkzexample}[latex=7cm,small]
\begin{tikzpicture}[scale=.5]
 \tkzDefPoints{0/0/O,-5/0/A,2/-2/B,0/5/D}
 \tkzInterLC[common=A](B,A)(O,D)
 \tkzGetFirstPoint{C}
 \tkzDrawPoints(O,A,B)
 \tkzDrawCircle(O,A)
 \tkzDrawLine(A,C)
 \tkzDrawPoint(C)
 \tkzLabelPoints(A,B,C)
\end{tikzpicture}
\end{tkzexample}


\subsubsection{Line-circle intersection: order of the points}
The idea is to compare the angles formed by the first point of the line, a resulting point and the center of the circle. The first point is the one that gives the smallest angle.

As you can see, $\widehat{BCO} < \widehat{BEO}$; in fact, $\widehat{BEO}$ is counterclockwise.

\begin{tkzexample}[latex=6cm,small]
\begin{tikzpicture}[scale=.5]
  \tkzDefPoints{0/0/O,5/1/A,2/2/B,3/1/D}
  \tkzInterLC[common=A](B,D)(O,A) \tkzGetPoints{C}{E}
  \tkzDrawPoints(O,A,B,D)
  \tkzDrawCircle(O,A) \tkzDrawLine(E,C)
  \tkzDrawSegments[dashed](B,O O,C)
  \tkzMarkAngle[->,size=1.5](B,C,O)
  \tkzDrawSegments[dashed](O,E)
  \tkzMarkAngle[->,size=1.5](B,E,O)
  \tkzDrawPoints(C,E)
  \tkzLabelPoints[above](O,E)
  \tkzLabelPoints[right](A,B,C,D)
\end{tikzpicture}
\end{tkzexample}

\subsubsection{Example with \tkzcname{foreach}}
\begin{tkzexample}[latex=7cm,small]
\begin{tikzpicture}[scale=3,rotate=180]
\tkzDefPoint(0,1){J} 
\tkzDefPoint(0,0){O}
\foreach \i in {0,-5,-10,...,-90}{
 \tkzDefPoint({2.5*cos(\i*pi/180)},%
   {1+2.5*sin(\i*pi/180)}){P}
 \tkzInterLC[R](P,J)(O,1)\tkzGetPoints{N}{M}
 \tkzDrawSegment[color=orange](J,N)
 \tkzDrawPoints[red](N)} 
\foreach \i in {-90,-95,...,-175,-180}{
 \tkzDefPoint({2.5*cos(\i*pi/180)},%
   {1+2.5*sin(\i*pi/180)}){P} 
 \tkzInterLC[R](P,J)(O,1)\tkzGetPoints{N}{M}
 \tkzDrawSegment[color=orange](J,M)
 \tkzDrawPoints[red](M)}   
\end{tikzpicture} 
\end{tkzexample}

\subsubsection{Line-circle intersection with option \tkzname{near}}
$D$ is the point closest to $b$.

\begin{tkzexample}[vbox,small]
  \begin{tikzpicture}
    \tkzDefPoints{0/0/A,12/0/C}
    \tkzDefGoldenRatio(A,C)                          \tkzGetPoint{B}
    \tkzDefMidPoint(A,C)                             \tkzGetPoint{O}
    \tkzDefMidPoint(A,B)                             \tkzGetPoint{O_1}
    \tkzDefMidPoint(B,C)                             \tkzGetPoint{O_2}
    \tkzDefPointBy[rotation=center O_2 angle 90](C)  \tkzGetPoint{P}
    \tkzDefPointBy[rotation=center O_1 angle 90](B)  \tkzGetPoint{Q}
    \tkzDefPointBy[rotation=center B angle 90](C)    \tkzGetPoint{b}
    \tkzInterLC[near](b,B)(O,A)                      \tkzGetFirstPoint{D}
    \tkzInterCC(D,B)(O,C)                            \tkzGetPoints{V}{U}
    \tkzDefPointBy[projection=onto U--V](O_1)        \tkzGetPoint{M}
    \tkzDefPointBy[projection=onto U--V](O_2)        \tkzGetPoint{N}  
    \tkzDrawPoints(A,B,C,O,O_1,O_2,D,U,V,M,N,b)
    \tkzDrawSemiCircles[teal](O,C O_1,B O_2,C)
    \tkzDrawSegments(A,C B,D U,V A,D C,D M,B B,N)
    \tkzDrawArc(D,U)(V)
    \tkzLabelPoints(A,B,C,O,O_1,O_2)
    \tkzLabelPoints[above](D,U,V,M,N)
  \end{tikzpicture}
\end{tkzexample}


\subsubsection{More complex example of a line-circle intersection}
Figure from \url{http://gogeometry.com/problem/p190_tangent_circle}

\begin{tkzexample}[latex=6.5cm,small]
\begin{tikzpicture}[scale=.75]
 \tkzDefPoint(0,0){A}  
 \tkzDefPoint(8,0){B}
 \tkzDefMidPoint(A,B)             \tkzGetPoint{O}
 \tkzDefMidPoint(O,B)             \tkzGetPoint{O'}
 \tkzDefLine[tangent from=A](O',B)\tkzGetFirstPoint{E}
 \tkzInterLC(A,E)(O,B)            \tkzGetFirstPoint{D}
 \tkzDefPointBy[projection=onto A--B](D)  
 \tkzGetPoint{F}
 \tkzDrawCircles(O,B O',B)
 \tkzDrawSegments(A,D A,B D,F) 
 \tkzDrawSegments[color=red,line width=1pt,
     opacity=.4](A,O F,B)
 \tkzDrawPoints(A,B,O,O',E,D) 
 \tkzMarkRightAngle(D,F,B)
 \tkzLabelPoints[below right](A,B,O,O',E,D) 
\end{tikzpicture}
\end{tkzexample}

\subsubsection{Circle defined by a center and a measure, and special cases}
Let's look at some special cases, such as straight lines tangent to the circle.

\begin{tkzexample}[latex=7cm,small]
\begin{tikzpicture}[scale=.5]
 \tkzDefPoint(0,8){A}      \tkzDefPoint(8,0){B}
 \tkzDefPoint(8,8){C}      \tkzDefPoint(4,4){D}
 \tkzDefPoint(2,4){E}      \tkzDefPoint(4,2){F}
 \tkzDefPoint(8,4){G}
 \tkzInterLC(A,C)(D,G)     \tkzGetPoints{I1}{I2}
 \tkzInterLC(B,C)(D,G)     \tkzGetPoints{J1}{J2}
 \tkzInterLC[near](A,B)(D,G)  \tkzGetPoints{K1}{K2}
 \tkzInterLC(E,F)(D,G)     \tkzGetPoints{E1}{E2}
 \tkzDrawCircle(D,G)
 \tkzDrawPoints[color=red](I1,J1,K1,K2,E1,E2)
 \tkzDrawLines(A,B B,C A,C I2,J2 E1,E2)
 \tkzDrawPoints[color=blue](A,...,F)
 \tkzDrawPoints[color=red](I2,J2)
 \tkzLabelPoints[left](B,D,E,F)
 \tkzLabelPoints[below left](A,C)
 \tkzLabelPoints[below=4pt](I1,K1,K2,E2)
 \tkzLabelPoints[left](J1,E1)
\end{tikzpicture}

\end{tkzexample}

\subsubsection{Calculation of radius}
 With \tkzname{pgfmath} and \tkzcname{pgfmathsetmacro}

The radius may be the result of a calculation, done before the intersection macro rather than inside it.
A length can be calculated in several ways. You can of course
 use the module \tkzname{pgfmath} and the macro \tkzcname{pgfmathsetmacro}. In some cases the results are not precise enough: the calculation $0.0002 \div 0.0001$ gives $1.98$ with pgfmath, while xfp gives $2$.

With \tkzname{xfp} and \tkzcname{fpeval}:

\begin{tkzexample}[latex=7cm,small]
  \begin{tikzpicture}
  \tkzDefPoint(2,2){A}
  \tkzDefPoint(5,4){B}
  \tkzDefPoint(4,4){O}
  \pgfmathsetmacro\tkzLen{\fpeval{0.0002/0.0001}}
 % or \edef\tkzLen{\fpeval{0.0002/0.0001}}
  \tkzInterLC[R](A,B)(O, \tkzLen)
  \tkzGetPoints{I}{J}
  \tkzDrawCircle(O,I)
  \tkzDrawPoints[color=blue](A,B)
  \tkzDrawPoints[color=red](I,J)
  \tkzDrawLine(I,J)
\end{tikzpicture}
  \end{tkzexample}


\subsubsection{Option \code{with nodes}}
\begin{tkzexample}[latex=8cm,small]
\begin{tikzpicture}[scale=.75]
\tkzDefPoints{0/0/A,4/0/B,1/1/D,2/0/E}
\tkzDefTriangle[equilateral](A,B)
\tkzGetPoint{C}
\tkzInterLC[with nodes](D,E)(C,A,B)
\tkzGetPoints{F}{G}
\tkzDrawCircle(C,A)
\tkzDrawPolygon(A,B,C)
\tkzDrawPoints(A,...,G)
\tkzDrawLine(F,G)
\end{tikzpicture}
\end{tkzexample}

\subsection{Intersection of two circles  \tkzcname{tkzInterCC}}

The most frequent case is that of two circles defined by their center and a point, but as before the option \tkzname{R} allows you to give the radius as a measure.

\begin{NewMacroBox}{tkzInterCC}{\oarg{options}\parg{$O,A$}\parg{$O',A'$} or \parg{$O,r$}\parg{$O',r'$} or   \parg{$O,A,B$} \parg{$O',C,D$}}%
\begin{tabular}{lll}%
options       & default & definition                         \\
\midrule
\TOline{N}   {N}    {$OA$ and $O'A'$ are radii, $O$ and $O'$ are the centers.}
\TOline{R}   {N}    {$r$ and $r'$ are dimensions and measure the radii.}
\TOline{with nodes} {N}  {in (A,A,C)(C,B,F) AC and BF give the radii. }
\TOline{common=pt}  {}   {pt is common point; tkzFirstPoint gives the other point.}
\bottomrule
\end{tabular}

\medskip
This macro defines the intersection points $I$ and $J$ of the two circles with centers $O$ and $O'$. If the two circles have no common point, the macro ends with an error that is not handled. If the centers are $O$ and $O'$ and the intersections are $A$ and $B$, the angles $\widehat{O,A,O'}$ and $\widehat{O,B,O'}$ have opposite orientations; \tkzname{tkzFirstPoint} is assigned to the point that forms the \code{clockwise} angle.
\end{NewMacroBox}

\begin{NewMacroBox}{tkzTestInterCC}{\parg{$O,A$}\parg{$O',B$}}%
The arguments are two pairs, which define two circles (each with a center and a point of the circle). If the two circles intersect, the test \tkzcname{iftkzFlagCC} is true.
\end{NewMacroBox}

\subsubsection{Testing a circle-circle intersection: \tkzcname{tkzTestInterCC}}

\begin{tkzexample}[latex=7cm,small]
\begin{tikzpicture}[scale=.75]
  \tkzDefPoints{% x   y   name
                   0    /0    /A,
                   2    /0    /B,
                   4    /0    /I,
                   1    /0    /P}
\tkzDrawCircle(A,B)
\foreach \i in {1,...,3}{%
   \coordinate  (P) at (\i,0);
\tkzDrawCircle[new](I,P)
   \tkzTestInterCC(A,B)(I,P)
    \iftkzFlagCC
    \tkzInterCC(A,B)(I,P)  \tkzGetPoints{a}{b}
    \tkzDrawPoints(a,b)
    \fi}
  \end{tikzpicture}
\end{tkzexample}

\subsubsection{Circle-circle intersection with a \tkzname{common} point}

\begin{tkzexample}[latex=7cm,small]
\begin{tikzpicture}[scale=.5]
  \tkzDefPoints{0/0/O,5/-1/A,2/2/B}
  \tkzDrawPoints(O,A,B)
  \tkzDrawCircles(O,B A,B)
  \tkzInterCC[common=B](O,B)(A,B)
  \tkzGetFirstPoint{C}
  \tkzDrawPoint(C)
  \tkzLabelPoints[above](O,A,B,C)
\end{tikzpicture}
\end{tkzexample}

\subsubsection{Circle-circle intersection: order of the points}
The idea is to compare the angles formed by the first center, a resulting point and the center of the second circle. The first point is the one that gives the smallest angle.

As you can see, $\widehat{ODB} < \widehat{OBE}$.

\begin{tkzexample}[latex=6cm,small]
\begin{tikzpicture}[scale=.5]
  \pgfkeys{/pgf/number format/.cd,fixed relative,
     precision=4}
  \tkzDefPoints{0/0/O,5/-1/A,2/2/B,2/-1/C}
  \tkzDrawPoints(O,A,B)
  \tkzDrawCircles(O,A B,C)
  \tkzInterCC(O,A)(B,C)\tkzGetPoints{D}{E}
  \tkzDrawPoints(C,D,E)
  \tkzLabelPoints(O,A,B,C)
  \tkzLabelPoints[above](D,E) 
  \tkzDrawSegments[cyan](D,O D,B)
  \tkzMarkAngle[red,->,size=1.5](O,D,B)
  \tkzFindAngle(O,D,B)   \tkzGetAngle{an}
  \tkzLabelAngle(O,D,B){$ \pgfmathprintnumber{\an}$}
  \tkzDrawSegments[cyan](E,O E,B)
  \tkzMarkAngle[red,->,size=1.5](O,E,B)  
  \tkzFindAngle(O,E,B)   \tkzGetAngle{an}
  \tkzLabelAngle(O,E,B){$ \pgfmathprintnumber{\an}$}  
\end{tikzpicture}
\end{tkzexample}

\subsubsection{Construction of an equilateral triangle}
$\widehat{A,C,B}$ is a clockwise angle.
\begin{tkzexample}[latex=7cm,small]
\begin{tikzpicture}[trim left=-1cm,scale=.5]
 \tkzDefPoint(1,1){A}
 \tkzDefPoint(5,1){B}
 \tkzInterCC(A,B)(B,A)\tkzGetPoints{C}{D}
 \tkzDrawPoint[color=black](C)
 \tkzDrawCircles(A,B B,A)
 \tkzCompass[color=red](A,C)
 \tkzCompass[color=red](B,C)
 \tkzDrawPolygon(A,B,C)
 \tkzMarkSegments[mark=s|](A,C B,C)
 \tkzLabelPoints[](A,B)
 \tkzLabelPoint[above](C){$C$}
\end{tikzpicture}
\end{tkzexample}


\subsubsection{Segment trisection}
 The idea here is to divide a segment into three equal parts with a ruler and a compass.

\begin{tkzexample}[latex=7cm,small]
\begin{tikzpicture}[scale=.6]
 \tkzDefPoint(0,0){A}
 \tkzDefPoint(3,2){B}
 \tkzInterCC(A,B)(B,A)          \tkzGetSecondPoint{D}
 \tkzInterCC(D,B)(B,A)          \tkzGetPoints{A}{C}
 \tkzInterCC(D,B)(A,B)          \tkzGetPoints{E}{B}
 \tkzInterLC[common=D](C,D)(E,D)\tkzGetFirstPoint{F}
 \tkzInterLL(A,F)(B,C)          \tkzGetPoint{O}
 \tkzInterLL(O,D)(A,B)          \tkzGetPoint{H}
 \tkzInterLL(O,E)(A,B)          \tkzGetPoint{G}
 \tkzDrawCircles(D,E A,B B,A E,A)
 \tkzDrawSegments[](O,F O,B O,D O,E)
 \tkzDrawPoints(A,...,H)
 \tkzDrawSegments(A,B B,D A,D A,E E,F C,F B,C)
 \tkzMarkSegments[mark=s|](A,G G,H H,B)
\end{tikzpicture}
\end{tkzexample}

\subsubsection{With the option \code{with nodes}}
\begin{tkzexample}[latex=6cm,small]
\begin{tikzpicture}[scale=.5]
 \tkzDefPoints{0/0/A,0/5/B,5/0/C}
 \tkzDefPoint(54:5){F}
 \tkzInterCC[with nodes](A,A,C)(C,B,F)
 \tkzGetPoints{a}{e}
 \tkzInterCC(A,C)(a,e) \tkzGetFirstPoint{b}
 \tkzInterCC(A,C)(b,a) \tkzGetFirstPoint{c}
 \tkzInterCC(A,C)(c,b) \tkzGetFirstPoint{d}
 \tkzDrawCircle[new](A,C)
 \tkzDrawPoints(a,b,c,d,e)
 \tkzDrawPolygon(a,b,c,d,e)
 \foreach \vertex/\num in {a/36,b/108,c/180,
                          d/252,e/324}{%
 \tkzDrawPoint(\vertex)
 \tkzLabelPoint[label=\num:$\vertex$](\vertex){}
 \tkzDrawSegment(A,\vertex)
 }
\end{tikzpicture}
\end{tkzexample}

\subsubsection{Mix of intersections}
\begin{tkzexample}[latex=8cm,small]
\begin{tikzpicture}[scale = .7]
  \tkzDefPoint(2,2){A}
  \tkzDefPoint(0,0){B}
  \tkzDefPoint(-2,2){C}
  \tkzDefPoint(0,4){D}
  \tkzDefPoint(4,2){E}
  \tkzCircumCenter(A,B,C)\tkzGetPoint{O}
  \tkzInterCC[R](O,2)(D,2)\tkzGetPoints{M1}{M2}
  \tkzInterCC(O,A)(D,O) \tkzGetPoints{1}{2}
  \tkzInterLC(A,E)(B,M1)\tkzGetSecondPoint{M3}
  \tkzInterLC(O,C)(M3,D)\tkzGetSecondPoint{L}
  \tkzDrawSegments(C,L)
  \tkzDrawPoints(A,B,C,D,E,M1,M2,M3,O,L)
  \tkzDrawSegments(O,E)
  \tkzDrawSegments[new](C,A D,B)
  \tkzDrawPoint(O)
  \tkzDrawCircles[new](M3,D B,M2 D,O)
  \tkzDrawCircle(O,A)
  \tkzLabelPoints[below right](A,B,C,D,E,M1,M2,M3,O,L)
\end{tikzpicture}
\end{tkzexample}


\subsubsection{Altshiller-Court's theorem}
  The two lines joining the points of intersection of two orthogonal circles to a point of one of the circles meet the other circle in two diametrically opposite points (Altshiller-Court, p.~176).


\begin{tkzexample}[vbox,small]
\begin{tikzpicture}
  \tkzDefPoints{0/0/P,5/0/Q,3/2/I}
  \tkzDefCircle[orthogonal from=P](Q,I) 
  \tkzGetFirstPoint{E}
  \tkzDrawCircles(P,E Q,E)
  \tkzInterCC[common=E](P,E)(Q,E) \tkzGetFirstPoint{F}
  \tkzDefPointOnCircle[through =  center P angle 80 point E] 
  \tkzGetPoint{A}
  \tkzInterLC[common=E](A,E)(Q,E)  \tkzGetFirstPoint{C}
  \tkzInterLL(A,F)(C,Q)  \tkzGetPoint{D}
  \tkzDrawLines[add=0 and 1](P,Q)
  \tkzDrawLines[add=0 and 2](A,E)
  \tkzDrawSegments(P,E E,F F,C A,F C,D)
  \tkzDrawPoints(P,Q,E,F,A,C,D)
  \tkzLabelPoints(P,Q,F)
  \tkzLabelPoints[above](E,A)
  \tkzLabelPoints[left](D)
  \tkzLabelPoints[above right](C)
\end{tikzpicture}
\end{tkzexample}


\endinput
\input{TKZdoc-euclide-angles.tex}

\part{Drawing, filling and clipping}
\input{TKZdoc-euclide-drawing.tex}
\input{TKZdoc-euclide-filling.tex}
\input{TKZdoc-euclide-clipping.tex}

\part{Marking and labelling}
\input{TKZdoc-euclide-marking.tex}
\input{TKZdoc-euclide-labelling.tex}

\part{Complements}
\input{TKZdoc-euclide-compass.tex}
\input{TKZdoc-euclide-show.tex}
\input{TKZdoc-euclide-rapporteur.tex}
\input{TKZdoc-euclide-tools.tex}

\part{Working with styles}
\input{TKZdoc-euclide-styles.tex}

\part{Examples}
\input{TKZdoc-euclide-others.tex}
\input{TKZdoc-euclide-examples.tex}

\part{FAQ}
\section{FAQ}

\subsection{Most common errors}
These are the mistakes I make most often myself, since the syntax has changed several times. This section will be expanded.

\begin{itemize}\setlength{\itemsep}{10pt}
  \item Forgetting the ``s'' of the macros that draw several objects, such as \tkzcname{tkzDrawSegments} or \tkzcname{tkzDrawCircles}: \tkzcname{tkzDrawPoint(A,B)} is wrong, you need \tkzcname{tkzDrawPoints(A,B)}.

  \item Since version 4, the unit is always the cm, so you must not write a unit: \tkzcname{tkzDefCircle[R](O,3cm)} becomes \tkzcname{tkzDefCircle[R](O,3)}. The usual \tkzname{TikZ} options keep their units, for example \tkzname{below right = 12pt}; on the other hand, you write \tkzname{size=1.2} to set the radius of the arc in \tkzcname{tkzMarkAngle}.

  \item Parentheses instead of braces. A point that is created has its name between braces, while a point that is used, as an option or as an argument, has its name between parentheses: write \tkzcname{tkzGetPoint\{A\}}, not \tkzcname{tkzGetPoint(A)}.

  \item \tkzcname{tkzGetPoint\{A\}} instead of \tkzcname{tkzGetFirstPoint\{A\}}. When a macro gives two points, either retrieve both with \tkzcname{tkzGetPoints\{A\}\{B\}}, or only one of them with \tkzcname{tkzGetFirstPoint\{A\}} or \tkzcname{tkzGetSecondPoint\{A\}}. These points can also be used with the references \tkzname{tkzFirstPointResult} and \tkzname{tkzSecondPointResult}. Some macros also give a third point, \tkzname{tkzPointResult}.

  \item Retrieve a result immediately: \tkzcname{tkzGetPoint}, \tkzcname{tkzGetPoints} or \tkzcname{tkzGetLength} must follow the macro that computes the result, because the next calculation replaces it.

  \item The order of the intersection points of a line and a circle, or of two circles, has sometimes changed between versions; this leads either to a wrong figure or to an error. The options \tkzname{common} and \tkzname{near} let you choose the point you want.

  \item Mixing options and arguments: all the macros that use a circle need to know its radius. If the radius is given by a measure, the option includes \tkzname{R}.

  \item Angles are given in degrees, more rarely in radians.

  \item A list of points must not be cut by a \code{\%} at the end of a line: |\tkzDrawSegments(A,B%| followed by |C,D)| gives the point |BC,D|, because \code{\%} also removes the space between the two points. Simply go to the next line without \code{\%}.

  \item The syntax of version 3 for special circles, for example |\tkzDrawCircle[euler](A,B,C)|, no longer exists: define the circle with \tkzcname{tkzDefCircle[euler](A,B,C)}, retrieve its center and a point with \tkzcname{tkzGetPoints\{N\}\{P\}}, then draw it with \tkzcname{tkzDrawCircle(N,P)}. Since version 5.14, the old syntax gives an error message.

  \item If an error occurs in a calculation when passing arguments, it is better to do the calculation before calling the macro.

  \item Do not mix the syntaxes of \tkzNamePack{pgfmath} and \tkzNamePack{xfp}. I have often chosen \tkzNamePack{xfp}; if you prefer \tkzNamePack{pgfmath}, do your calculations before passing the arguments.

  \item The error ``dimension too large'' occurs in some cases, for example with \tkzcname{tkzDrawArc} or \tkzcname{tkzMarkAngle} and a very small angle: it comes from the library \NameLib{decorations} when a mark is placed on the arc, even when the mark is absent. One way to avoid it is the option \tkzname{veclen} (see section \ref{opt-veclen}): in the scope where it is used, the function \code{veclen} is computed with \tkzNamePack{xfp}. Do not use intersection macros in this scope.

  \item Since version 5.14, \tkzcname{tkzSetUpCompass}\code{[delta=...]} no longer extends the arcs drawn with \tkzcname{tkzDrawArc}: use \tkzcname{tkzSetUpArc}\code{[delta=...]} for them.
\end{itemize}
\endinput


\part{News, compatibility and options}
\section{News and compatibility}\label{news}

\subsection{With version 5.15}
\begin{itemize}
  \item \tkzcname{tkzInterLC} with a line passing through the center of the circle: the order of the two points no longer depends on rounding. \tkzname{tkzSecondPoint} is the point located in the direction of the line (from its first point to its second point); the results are the same with and without the option \tkzname{lua}. See section \ref{interLC-center}.
\end{itemize}

\subsection{With version 5.14}
\begingroup\sloppy

This version only corrects bugs; the syntax does not change and no macro has been added.

\begin{itemize}
  \item \tkzcname{tkzDefCircle} with the options \tkzname{orthogonal from} or \tkzname{orthogonal through} no longer disturbs the following calls of \tkzcname{tkzDefCircle}.
  \item The options of the following macros now only apply to the current call: \tkzcname{tkzShowLine}, \tkzcname{tkzDefSimilitudeCenter}, \tkzcname{tkzCalcLength}, \tkzcname{tkzGrid}, \tkzcname{tkzClip}, \tkzcname{tkzDefEquiPoints}, \tkzcname{tkzLegend}.
  \item \tkzcname{tkzIsLinear} gave a wrong answer (always true); \tkzcname{tkzIsLinear} and \tkzcname{tkzIsOrtho} now use a tolerance that depends on the lengths.
  \item The result of \tkzcname{tkzPowerCircle} no longer changes after another calculation.
  \item The inversion of a circle (\tkzcname{tkzDefCircleBy}) gives a line only if the circle passes through the center of inversion.
  \item \tkzcname{tkzDefCircle}\code{[spieker]} gives the radius in cm, like the other options.
  \item The Feuerbach point and the Feuerbach triangle, obtained with the option \tkzname{feuerbach} of \tkzcname{tkzDefTriangleCenter} and \tkzcname{tkzDefSpcTriangle}, are now exact.
  \item \tkzcname{tkzDrawPolySeg}: the arrow at the start of the path has the right direction.
  \item \tkzcname{tkzCompass} accepts a decimal value for \tkzname{delta}; \tkzcname{tkzProtractor}\code{[return]} works; \tkzcname{tkzLabelRegPolygon} uses its option \tkzname{sep}; \tkzcname{tkzDrawX} uses its option \tkzname{step}.
  \item The old syntax of version 3, for example |\tkzDrawCircle[euler](A,B,C)|, now gives an error message instead of an endless compilation.
  \item The aliases \tkzcname{tkzDefNinePointsCircle}, \tkzcname{tkzFeuerBachCircle} and \tkzcname{tkzIncentralTriangle} work again.
\end{itemize}

Two changes may modify your figures:
\begin{itemize}
  \item \tkzcname{tkzSetUpCompass}\code{[delta=...]} only acts on the compass marks. To extend the arcs drawn by \tkzcname{tkzDrawArc}, \tkzcname{tkzMarkArc} and \tkzcname{tkzLabelArc}, use \tkzcname{tkzSetUpArc}\code{[delta=...]}. Before this version, the same value was used for both.
  \item The \TIKZ\ style \tkzname{help lines} is no longer modified by the package.
\end{itemize}

\endgroup

\subsection{With version 5.10}
\begin{itemize}
  \item Added a French documentation.
  \item Added the \code{mini} option. You can use this option with the \tkzNamePack{tkz-elements} package. Only the modules required for drawing will be loaded. This option is currently only useful if you are using \tkzNamePack{tkz-elements}.
\end{itemize}

\subsection{From version 3 to version 5}

\tkzname{\tkznameofpack} 4 was a simplified version of its predecessor, independent of \tkzname{tkz-base}. Most macros of \tkzname{tkz-euclide} 3 have been retained, but the special circles are now defined with \tkzcname{tkzDefCircle} and then drawn with \tkzcname{tkzDrawCircle}. The unit is now the \tkzname{cm}. If you need some macros from \tkzname{tkz-base}, you may need to use \tkzcname{tkzInit}.

Version 5 does not change anything for the user. With the option \code{lua} (|\usepackage[lua]{tkz-euclide}|), the calculations are done by Lua: you must then compile with LuaLaTeX. See section \ref{calc_with_lua}.



\endinput
\input{TKZdoc-euclide-lua.tex}

\clearpage\newpage
\small\printindex
\end{document}